\section{Test 3: síntesi de la unitat}
\iniciTest{r3}{c,d,b,a,c,b,b,a,b,b,c,c}

\begin{enumerate}
\pregunta Quin conjunt numèric conté exactament els nombres que es poden escriure com a quocient de dos enters, amb denominador no nul?
\begin{enumerate}
\opcio{a}{$\mathbb{N}$}
\opcio{b}{$\mathbb{Z}$}
\opcio{c}{$\mathbb{Q}$}
\opcio{d}{$\mathbb{R}\setminus\mathbb{Q}$}
\end{enumerate}
\feedback


\pregunta Quina afirmació descriu correctament els nombres irracionals?
\begin{enumerate}
\opcio{a}{Són decimals finits.}
\opcio{b}{Són decimals periòdics.}
\opcio{c}{No es poden representar sobre la recta real.}
\opcio{d}{No es poden expressar com a fracció de dos enters.}
\end{enumerate}
\feedback


\pregunta Si $A=(-2,5]$ i $B=[1,7)$, quin és $A\cap B$?
\begin{enumerate}
\opcio{a}{$(-2,7)$}
\opcio{b}{$[1,5]$}
\opcio{c}{$(1,5)$}
\opcio{d}{$[-2,7]$}
\end{enumerate}
\feedback


\pregunta Quin és el conjunt solució de $|x-3|<2$?
\begin{enumerate}
\opcio{a}{$(1,5)$}
\opcio{b}{$(-5,-1)$}
\opcio{c}{$(3,5)$}
\opcio{d}{$(-2,2)$}
\end{enumerate}
\feedback


\pregunta Quina desigualtat és sempre certa si $0<a<b$?
\begin{enumerate}
\opcio{a}{$a^2>b^2$}
\opcio{b}{$-a<-b$}
\opcio{c}{$\frac{1}{a}>\frac{1}{b}$}
\opcio{d}{$\frac{1}{a}<\frac{1}{b}$}
\end{enumerate}
\feedback


\pregunta Quin és el suprem de l’interval $(2,7)$?
\begin{enumerate}
\opcio{a}{$2$}
\opcio{b}{$7$}
\opcio{c}{No existeix perquè $7$ no pertany a l’interval.}
\opcio{d}{$\frac{9}{2}$}
\end{enumerate}
\feedback


\pregunta Quin és el resultat de simplificar $\sqrt{50}-\sqrt{8}$?
\begin{enumerate}
\opcio{a}{$\sqrt{42}$}
\opcio{b}{$3\sqrt{2}$}
\opcio{c}{$5\sqrt{2}$}
\opcio{d}{$7\sqrt{2}$}
\end{enumerate}
\feedback


\pregunta Quina expressió és equivalent a $a^{3/2}$, si $a>0$?
\begin{enumerate}
\opcio{a}{$\sqrt{a^3}$}
\opcio{b}{$\sqrt[3]{a^2}$}
\opcio{c}{$a\sqrt[3]{a}$}
\opcio{d}{$\frac{1}{\sqrt{a^3}}$}
\end{enumerate}
\feedback


\pregunta En racionalitzar $\frac{1}{\sqrt{3}}$, s’obté:
\begin{enumerate}
\opcio{a}{$\sqrt{3}$}
\opcio{b}{$\frac{\sqrt{3}}{3}$}
\opcio{c}{$\frac{1}{3}$}
\opcio{d}{$3\sqrt{3}$}
\end{enumerate}
\feedback


\pregunta Quin valor té $\sqrt[3]{27}$?
\begin{enumerate}
\opcio{a}{$2$}
\opcio{b}{$3$}
\opcio{c}{$4$}
\opcio{d}{$9$}
\end{enumerate}
\feedback


\pregunta Si $10^{-2}=0.01$, quin ordre de magnitud expressa aquesta potència?
\begin{enumerate}
\opcio{a}{Centenes.}
\opcio{b}{Desenes.}
\opcio{c}{Centèsimes.}
\opcio{d}{Mil·lèsimes.}
\end{enumerate}
\feedback


\pregunta Quina afirmació resumeix millor la recta real?
\begin{enumerate}
\opcio{a}{Només conté nombres racionals.}
\opcio{b}{Té buits corresponents als irracionals.}
\opcio{c}{Cada punt correspon a un únic nombre real i cada nombre real a un punt.}
\opcio{d}{Només serveix per representar enters.}
\end{enumerate}
\feedback

\end{enumerate}